%% Lebenslauf - Ana Lorena Ortiz Loyola
%% Zugeschnitten auf: Fraunhofer SIT (ATHENE) - Studentische Hilfskraft zur Unterstuetzung der ATHENE Geschaeftsstelle, Darmstadt
%% Deutsche Version von main_fraunhofer_athene.tex.
%%
%% Kompilieren mit: lualatex main_fraunhofer_athene_de.tex
%% (pdflatex schlägt auf modernem MiKTeX häufig mit fontawesome5-Fehlern fehl.)

\documentclass[11pt,a4paper,sans]{moderncv}
\moderncvstyle{banking}
\moderncvcolor{blue}

\usepackage{needspace}

% Erzwingt, dass Vor- und Nachname sowie die Sektionsüberschriften in moderncv
% Blau (color1) dargestellt werden. Der banking-Standardstil rendert diese auf
% lualatex+MiKTeX sonst schwarz, was inkonsistent mit dem restlichen blauen
% Akzentschema wirkt.
\renewcommand*{\firstnamestyle}[1]{{\fontsize{34}{36}\bfseries\upshape\color{color1}#1}}
\renewcommand*{\lastnamestyle}[1]{{\fontsize{34}{36}\bfseries\upshape\color{color1}#1}}
\renewcommand*{\sectionstyle}[1]{{\sectionfont\color{color1}#1}}

\usepackage[utf8]{inputenc}
\usepackage{hyperref}
\hypersetup{
    colorlinks=true,
    linkcolor=blue,
    filecolor=magenta,
    urlcolor=blue,
    pdftitle={Ana Lorena Ortiz Loyola - Lebenslauf},
    pdfpagemode=FullScreen,
}
\usepackage[scale=0.80]{geometry}
\usepackage{import}

% Persönliche Daten
\name{Ana Lorena}{Ortiz Loyola}
\address{Frankfurt am Main, Deutschland}{}{}
\phone[mobile]{+49 176 62145797}
\email{a01750283@tec.mx}
\extrainfo{\href{https://linkedin.com/in/lorenasolana}{LinkedIn}}

\begin{document}

\makecvtitle

% ============================================================
%     PROFILTEXT
% ============================================================

\vspace{6pt}
\small{Data-Science-Studentin (Tecnológico de Monterrey, Abschluss voraussichtlich 2027) mit einer strukturierten, detailgenauen Arbeitsweise, die sich in mehreren Forschungs- und Datenqualitätsprojekten gezeigt hat, sicheren MS-Office- und Excel-Kenntnissen sowie praktischer Erfahrung in der Koordination von Dokumentation und Stakeholder-Kommunikation während eines Data-Engineering-Praktikums bei der Deutschen Bundesbank (AnaCredit). Zuverlässig, organisiert und geübt im eigenständigen Arbeiten an administrativen und koordinierenden Aufgaben.}

% ============================================================
%     KERNKOMPETENZEN
% ============================================================

\section{Kernkompetenzen}
\vspace{1pt}
\begin{itemize}

\item \textbf{Büroorganisation}: Excel, Word, PowerPoint, strukturierte Dokumentation und Berichterstattung

\item \textbf{Stakeholder-Kommunikation}: Koordiniertes automatisiertes Reporting und direkte Kommunikation mit externen Ansprechpartnern (Deutsche Bundesbank AnaCredit)

\item \textbf{Detailgenaue Arbeitsweise}: Nachvollziehbare, vollständig dokumentierte Arbeit über alle technischen und organisatorischen Projekte

\item \textbf{Daten \& Programmierung}: Python, SQL, statistische Analyse (ergänzende Fähigkeiten über die Kernaufgaben hinaus)

\item \textbf{Sprachen}: Deutsch (C1), Englisch (C1)

\end{itemize}

% ============================================================
%     BERUFSERFAHRUNG
% ============================================================

\section{Berufserfahrung}
\vspace{3pt}
\begin{itemize}

\needspace{5\baselineskip}
\item{\cventry{04/2026--07/2026}{Data Engineer / Praktikantin}{Deutsche Bundesbank (AnaCredit)}{Frankfurt am Main, Deutschland}{}{\vspace{1pt}
\begin{itemize}
    \item Entwickelte automatisierte Migrationstests und Dokumentation, um identische Geschäftslogik nach der Datenmigration sicherzustellen
    \item Koordinierte automatisiertes Reporting und direkte Kommunikation mit meldenden Banken zu festgestellten Datenqualitätsproblemen
\end{itemize}}}

\end{itemize}

% ============================================================
%     PROJEKTE
% ============================================================

\needspace{8\baselineskip}
\section{Projekte}
\vspace{1pt}
\begin{itemize}

\needspace{5\baselineskip}
\item{\cventry{05/2026}{Banking-Datenpipeline -- Datenbereinigung \& Prüfprotokoll}{Persönliches Projekt (GitHub), umgesetzt mit Claude Code}{}{}{\vspace{1pt}
\begin{itemize}
    \item Entwickelte ein vollständig nachvollziehbares Datenbereinigungsprotokoll, das jede Änderung (Tabelle, Spalte, alter/neuer Wert, Aktion, Zeitstempel) für simulierte Banktransaktionen festhält
    \item Erstellte Qualitätsvisualisierungen und einen strukturierten Ausnahmebericht
\end{itemize}}}

\vspace{3pt}

\item{\cventry{10/2024--12/2024}{Predictive Modeling -- Optimierung im Gesundheitswesen}{Tecnológico de Monterrey}{}{}{\vspace{1pt}
\begin{itemize}
    \item Analysierte über 25 Jahre nationaler Gesundheitsdaten und dokumentierte Ergebnisse in strukturierten Visualisierungen
\end{itemize}}}

\vspace{3pt}

\item{\cventry{10/2024--02/2025}{Forschung zu gitterbasierter Kryptografie}{Tecnológico de Monterrey}{}{}{\vspace{1pt}
\begin{itemize}
    \item Implementierte und analysierte gitterbasierte kryptografische Algorithmen (LWE, NTRU) in Python; Beitrag zu einer wissenschaftlichen Publikation
\end{itemize}}}

\end{itemize}

% ============================================================
%     AUSBILDUNG
% ============================================================

\section{Ausbildung}
\vspace{1pt}
\begin{itemize}

\item{\cventry{2022--2027}{BSc in Data Science}{Tecnológico de Monterrey}{Mexiko}{}{\vspace{1pt}
Note: 1,3 (Stipendium für akademische Exzellenz; DAAD-KOSPIE-Stipendium). Schwerpunkte: Datenanalyse \& Statistik, Kryptografie \& Sicherheit.
}}

\vspace{3pt}

\item{\cventry{2025--2026}{Auslandssemester}{Universität Göttingen}{Deutschland}{}{}}

\end{itemize}

% ============================================================
%     SPRACHEN
% ============================================================

\section{Sprachen}
\vspace{1pt}
\begin{itemize}
\item Deutsch (C1), Englisch (C1, TOEFL iBT 105/120), Spanisch (Muttersprache).
\end{itemize}

% ============================================================
%     PUBLIKATIONEN
% ============================================================

\section{Publikationen}
\vspace{1pt}
\begin{itemize}
\item Beitrag zu einer wissenschaftlichen Publikation über gitterbasierte Kryptografie (Sicherheitsparameter- und Laufzeitanalyse von LWE und NTRU), gemeinsam mit Prof. A. F. De Abiega L'Eglisse, Tecnológico de Monterrey (Zitation ausstehend, Publikation in Vorbereitung).
\end{itemize}

% ============================================================
%     AUSZEICHNUNGEN
% ============================================================

\section{Auszeichnungen}
\vspace{1pt}
\begin{itemize}
\item{\textbf{Stipendium für akademische Exzellenz} -- Tecnológico de Monterrey.}
\item{\textbf{DAAD-KOSPIE-Stipendium}.}
\item{\textbf{Nationale Physik-Olympiade} -- Mexiko (2020).}
\item{\textbf{Lobende Erwähnung, Physik-Olympiade} (2021).}
\end{itemize}

% ============================================================
%     REFERENZEN
% ============================================================

\section{Referenzen}
\vspace{1pt}
\begin{itemize}
\item{Auf Anfrage erhältlich.}
\end{itemize}

\end{document}
